\chapter{How the control approach reached its present form}
\label{ch:development}

The system described in \cref{ch:method} is not the one the planning report
proposed, and the route between them is part of the result. This chapter
records the sequence, the alternatives that were considered at each branch,
and what each abandoned approach contributed to the one that shipped. The
divergence from the planning report is analysed separately in
\repo{thesis\_report/DIVERGENCE.md}.

The account below is no longer reconstructed. The prior programs were
recovered and are archived at \repo{thesis\_report/\_source/prior\_work/};
file dates give the calendar directly, running from 7 June to 27 July 2026.

\section{Phase 1: joint-space control in simulation}

The planning report proposed validating the master and slave mapping in Isaac
Lab or MuJoCo before touching hardware. MuJoCo~\cite{todorov2012mujoco} was chosen over Isaac
Lab~\cite{isaaclab}. It starts in under a second, its model format
is a single readable file, and for a kinematic mapping question its contact
solver is irrelevant, so the heavier alternative bought nothing.

The first mapping was the obvious one. The master reproduces the Gen3's joint
sequence, so each master potentiometer angle was sent to the corresponding
slave joint after zero-referencing. Joint one to joint one, and so on through
seven.

Three things were learned, and all three survived into the delivered system.

\paragraph{The mapping is dominated by convention handling, not by control.}
Almost all of the difficulty was in what a joint angle \emph{means}: which
direction is positive, where zero sits, whether the value wraps, and over what
interval. This is where the degree and radian conversion described in
\cref{ch:intro} was first written, and it is the reason it later became an
isolated, unit-tested module rather than an inline expression. The lesson
transferred directly: the arithmetic that converts between conventions
deserves more test coverage than the control law it feeds.

\paragraph{An unconstrained joint mapping will drive the arm into things.}
In simulation this is a visual annoyance. On the delivered platform, where the
arm is mounted beside a person's head, it is the entire safety problem. A
joint-space mapping has no representation of the wearer, so nothing in the
formulation can express the constraint that matters most. That observation is
what motivated the move to a framework with a collision model in it.

\paragraph{Sensor faults are the dominant failure mode, not control faults.}
Injecting a stuck or noisy channel showed the arm following it faithfully,
because a direct joint mapping has no notion of a reading being implausible.
The channel-health architecture in \cref{ch:method} is a direct response.

\subsection{What the prior programs actually did}

Five of the recovered programs import MuJoCo and none contains an inverse
kinematics call of any kind. The mapping is written literally as a joint-space
assignment: \verb|data.qpos[k1_qpos[j]] = k1_target[j]| in
\repo{two\_armsbackpack.py}, and the same pattern by way of degrees in
\repo{srl\_teleop.py}. Searching the whole set for \emph{inverse} or
\emph{jacobian} returns nothing. The phase was therefore joint-space
throughout, exactly as described, and the absence of any solver is what makes
the later move to a task-space command a change of kind rather than of degree.

Two facts recovered from the source correct the account given previously.

\paragraph{The prior system already drove the real arm.} \repo{srl\_teleop.py}
imports the manufacturer's Python API directly and calls
\verb|SendJointSpeedsCommand|, gated behind a \verb|USE_REAL| mode. It is the
\emph{same high-level velocity interface} the delivered system uses, arrived
at independently and about two months earlier. So the move to ROS~2 was not
driven by a need for vendor drivers, as an earlier draft of this chapter
supposed: the arm was already reachable. It was driven by the collision model,
the solver and the controller lifecycle, which is a narrower and more honest
claim.

\paragraph{The link lengths have a documented origin.} The seven values the
forward kinematics uses were not first written in Python. They appear in
\repo{basic\_control/kinematics.h} as C++ constants, with the comment
``physical link lengths (metres), shaft-center to shaft-center, measured
directly on the physical master arm'', alongside the joint pattern
\verb|J1 roll, J2 bend, ... J7 roll|.

That provenance settles what the CAD comparison in \cref{ch:results} can
claim. \emph{Shaft-centre to shaft-centre} is precisely the joint-module pitch
measured off the STEP assembly, so the two quantities are commensurable, and
their agreement to \SI{1.2}{\percent} is a genuine cross-check rather than a
coincidence of similar numbers. It also explains why the individual
\SI{43}{\milli\metre} and \SI{37}{\milli\metre} figures cannot be separated
from the CAD: neither the header nor the geometry distinguishes them, only
their sum.

\subsection{Why joint-space control was abandoned}

Joint-space correspondence has a genuine advantage that should be recorded
before the case against it: it fixes the redundant degree of freedom for free.
The Gen3 has seven joints for a six-dimensional task, so a task-space command
leaves one dimension unspecified and something has to choose it. Under a joint
mapping the operator chooses it implicitly, by the posture they put the master
in.

It was abandoned for two reasons that are both about the master rather than
the method.

The first is scale. The master's links sum to \SI{0.272}{\metre} against the
Gen3's \SI{0.902}{\metre} reach. Equal joint angles therefore produce tip
displacements differing by a factor of about three, so the operator's motion
and the robot's motion are related by a gain that varies with posture. Working
in task space makes that gain an explicit, tunable scalar instead.

The second is decisive. A joint mapping requires \emph{every} joint to be
measured, because each one contributes directly to the commanded
configuration. It has no graceful degradation at all: a single failed channel
sends a wrong angle to the corresponding slave joint, and there is no
reconstruction available because nothing else in the mapping observes that
joint. Given the argument in \cref{ch:intro} that a body-worn master will lose
channels, a formulation that cannot survive losing one was not viable.

\section{Phase 2: the move to ROS 2}

MuJoCo answers questions about dynamics. It does not supply a collision-aware
inverse kinematics service, a controller stack, or drivers for the actual
arms, and by the end of phase 1 all three were needed. The system moved to ROS
2 Jazzy~\cite{macenski2022ros2}.

Four capabilities came with the move, and each closed a problem left open by
phase 1.

\begin{itemize}
  \item \textbf{A collision model containing the wearer.} MoveIt checks
    candidate configurations against a scene, so the wearer's body can be part
    of the model. This is the capability a joint mapping could not express at
    all.
  \item \textbf{Inverse kinematics with redundancy resolution.} TRAC-IK solves
    the seven-joint problem with random restarts and a distance
    objective~\cite{beeson2015tracik}, which makes task-space commanding
    practical and lets the redundant degree of freedom be re-seeded on
    failure.
  \item \textbf{A controller stack with a hardware abstraction.} The same
    trajectory interface addresses a mock system and the real driver, so
    everything upstream is developed and tested unchanged and the substitution
    is one launch argument.
  \item \textbf{Vendor drivers.} The manufacturer supplies a ROS 2 driver, so
    the alternative was writing one.
\end{itemize}

The cost was real. ROS 2 introduced a discovery layer, a middleware
configuration, and a process lifecycle, and a substantial fraction of the
faults recorded in this project are in that layer rather than in the robot.
\Cref{ch:results} reports them, because they are the ones a successor will
meet first.

\subsection{What was kept from simulation}

Simulation did not stop being useful when hardware became available. It became
the place where anything that could not be attempted on a person is
attempted. Every workspace measurement in \cref{ch:workspace}, every fault
injection, and every clearance figure is simulated, and that is a deliberate
allocation rather than a fallback. The mock hardware interface is what makes
it possible to exercise the entire stack, including the trajectory controllers
and the emergency stop, without an arm attached.

The limit of that substitution is stated where it binds. The mock returns
identically zero for every force, so anything derived from a force reading is
structurally undefined rather than merely unmeasured.

\section{Phase 3: from Cartesian control to a reduced-degree-of-freedom
command}

Task-space control was not the end of the sequence. The first Cartesian
implementation computed the master's tip pose by forward kinematics from all
seven potentiometer angles and commanded the end effector to the scaled
displacement. That is the textbook formulation and it failed on this hardware.

The failure was directional rather than noisy. Decomposing the commanded
positions from five distinct gestures put \SI{93.4}{\percent} of the variance
on a single axis, and placed the pose commanded by raising the master
\emph{below} the pose commanded by lowering it. A mapping that inverts up and
down is not a mapping with a gain error.

The cause is that forward kinematics compounds every joint angle in the chain.
An error in joint three propagates through joints four to seven, so the tip
direction inherits the sum of seven absolute-angle calibration errors, and the
inertial mount calibration had already shown those errors to be large enough
to be physically impossible in places.

The response was to stop computing direction from the chain. Position is now
built in spherical form about the master's shoulder, taking each component
from whichever sensor observes it most directly: elevation from the gravity
vector, azimuth from the single shoulder-roll potentiometer, and radius from
the forward-kinematic magnitude, which depends only on the bend joints and is
insensitive to roll errors. \Cref{ch:method} gives the formulation and
\cref{ch:results} the measured consequences, including the one axis that this
does not fix.

\subsection{Alternatives considered at this branch}

\begin{description}
  \item[Recalibrate the potentiometers and keep full forward kinematics.]
    Rejected on cost and on durability. It requires an external angular
    reference per joint, and it has to be redone whenever a joint is
    reassembled. It also fixes nothing about channel loss, since forward
    kinematics still needs every joint.
  \item[Take azimuth from forward kinematics instead of from joint one.] Built
    and measured, then rejected. It made direction worse on both arms, moving
    the mean reversal error from \SI{12.6}{\degree} to \SI{33.3}{\degree} on
    the left arm and from \SI{26.2}{\degree} to \SI{39.5}{\degree} on the
    right. It fails for the same reason the full chain fails: it compounds
    four suspect angles where the single-joint estimate compounds one.
  \item[Integrate the gyroscope for azimuth.] Implemented, measured, and left
    disabled by default. It observes rotation about vertical directly and so
    avoids both error sources, and it halves the residual on the left arm. It
    is not enabled because one of the two sensors drifts at
    \SI{44}{\degree\per\minute} while stationary, and because the capture that
    would validate it needs a moving master and does not yet exist.
  \item[Add a second inertial sensor on the upper arm.] Not attempted, for the
    hardware reason in \cref{ch:master}: both available addresses are
    occupied. It remains the most promising route, and
    \cref{ch:results} quantifies what it would buy.
\end{description}

\section{Summary of the sequence}

\begin{table}[htbp]
  \centering
  \caption{Development sequence, with the reason each stage ended.}
  \label{tab:devpath}
  \begin{tabular}{@{}p{31mm}p{34mm}p{62mm}@{}}
    \toprule
    Stage & Mapping & Why it ended \\
    \midrule
    MuJoCo, phase 1 & joint to joint &
      No collision model, so the wearer could not be represented. No
      degradation path: one failed channel corrupts one slave joint
      irrecoverably. \\
    ROS 2, first Cartesian & tip pose by full forward kinematics &
      Direction collapsed. \SI{93.4}{\percent} of variance on one axis, and
      raising the master lowered the arm. Seven compounded calibration errors. \\
    ROS 2, delivered & spherical, per-component sensing &
      Current. Elevation and reach are correct; azimuth remains limited by the
      lever-arm coupling described in \cref{ch:results}. \\
    \bottomrule
  \end{tabular}
\end{table}

The through-line is that each stage failed because of what it assumed about
the \emph{sensors}, not because of what it assumed about the control. Joint
mapping assumed every channel works. Full forward kinematics assumed every
angle is accurate. The delivered mapping assumes only that gravity points down
and that the best-conditioned joints are identifiable at run time, which are
assumptions the hardware can actually support.
